\section{Experimental Evaluation}
\label{sec:experiments}
\subsection{Experimental Setup and Calibration}
To establish precise mathematical and geometric boundaries, all experiments are conducted within a mathematically controlled, synthetic 12-layer sequential PyTorch representation sandbox. This sandbox is modeled after a frozen Vision Transformer backbone (\texttt{vit\_tiny\_patch16\_224}, $D=192$) and represents four task domains: MNIST, FashionMNIST, CIFAR-10, and SVHN. While this controlled environment provides algebraic tractability over representation subspaces and linear covariate shifts, we recognize that synthetic Gaussian manifolds represent an idealized setting compared to real image embeddings (which exhibit topological overlaps and non-isotropic noise). In particular, the assumption of perfect subspace and class orthogonality in our sandbox may artificially amplify the \emph{Representational Sparsity Paradox} compared to smoother, correlated manifolds in real-world neural network representation spaces (such as pre-trained LLMs or Vision Transformers). In real backbones, class manifolds have substantial topological overlap, which would allow centroid-based methods to decay more gracefully while enabling our soft temperature ensembling to act as an even stronger spatial interpolator. We provide a detailed theoretical analysis of these architectural assumptions and sandbox simplifications in Section~\ref{sec:limitations}.
\begin{itemize}
    \item \textbf{Subspace Partitioning:} The 192-dimensional representation space is divided into four orthogonal 48-dimensional task subspaces, guaranteeing that task-specific features are initially orthogonal.
    \item \textbf{Class Orthogonality:} Within each 48-dimensional subspace, 10 orthogonal class prototypes of unit norm are generated via QR decomposition. Isotropic Gaussian noise is injected with calibrated scales matching domain complexity: MNIST (0.01), F-MNIST (0.05), CIFAR-10 (0.24), and SVHN (0.56).
    \item \textbf{Expert Training:} Independent expert models are trained with LoRA adapters (rank $r=8$) at Layers 4 to 12 and task-specific classification heads for 60 epochs using AdamW. The stand-alone expert accuracies serve as our \textbf{Expert Ceiling}:
    MNIST ($100.00\%$), FashionMNIST ($100.00\%$), CIFAR-10 ($90.88\%$), SVHN ($39.44\%$), yielding a Joint Mean ceiling of \textbf{$82.58 \pm 0.94\%$}. Specifically, the SVHN expert ceiling is calibrated to $39.44\%$ by design due to an intentionally high noise scale ($0.56$). This severe noise scale represents highly degraded edge sensors or low-light cameras, acting as an aggressive real-world stress-test for out-of-task noise rejection, routing stability, and cross-task noise infiltration under severe corruption. While this high noise skews the Joint Mean downwards for all models, it serves to critically magnify the differences in cross-task noise infiltration between methods. Under uniform or weaker routing, SVHN's high noise actively pollutes cleaner task spaces (like MNIST or CIFAR-10), degrading their accuracy. By showing that EER preserves near-perfect MNIST and FashionMNIST performance despite this extreme SVHN noise, we empirically demonstrate that direct prediction-entropy-based routing is exceptionally robust against such noise infiltration compared to centroid-based ensembling, where representation-space sparsity allows noisy activations to easily pollute task routing boundaries.
    \item \textbf{Overlapping Class Namespace and Evaluation Bias:} In both the synthetic sandbox and the real-world ResNet-18 experiments, all $K=4$ tasks have exactly $C=10$ classes, and their label spaces overlap within the integer namespace $\{0, \dots, 9\}$. Because sample-level accuracy is evaluated via a standard argmax match \texttt{(pred == y)}, incorrect routing can theoretically result in a ``false correct'' prediction. For example, if a CIFAR-10 sample with ground-truth class index 3 (``cat'') is routed incorrectly to the MNIST expert, and the MNIST expert predicts class index 3 (``3''), the evaluation script will register this prediction as correct. This overlapping namespace introduces a background chance probability of $\approx 10\%$ (representing a $1/C$ chance under a uniform prediction distribution) for misrouted samples, creating a slight optimistic bias in absolute accuracy scores across all ensembling models. While the well-separated and highly orthogonal features of the sandbox and ResNet-18 limit the empirical frequency of such cross-task prediction matches, they represent an inherent property of multi-task ensembling evaluations with overlapping class counts, which future benchmarks can mitigate by employing disjoint label spaces (e.g., $Y_k = [10k, 10k+9]$).
\end{itemize}

\textbf{SVHN Noise and 3-Task Joint Mean Analysis:} To justify the choice of this extreme noise scale for SVHN and quantify its impact on our ensembling evaluations, we conduct a brief ablation by calculating the Joint Mean across the three clean domains (MNIST, FashionMNIST, CIFAR-10) by excluding SVHN. Under EER, the clean 3-task Joint Mean is \textbf{88.13\%}, representing a remarkably small gap of only \textbf{-8.83\%} relative to the perfect-routing 3-task Expert Ceiling of \textbf{96.96\%}. Conversely, the 3-task heterogeneous SOTA baseline (SPS-ZCA) achieves only \textbf{82.61\%}, meaning EER outpaces the supervised baseline by \textbf{+5.52\%} absolute. Uniform Merging and EPL-OCA (Refined) yield clean joint means of \textbf{69.20\%} and \textbf{60.65\%}, respectively. This highlights that while SVHN's high noise scale (0.56) acts as an aggressive stress-test that skews the absolute 4-task Joint Mean downwards for all models, our primary proposed direct-routing paradigm (EER) remains exceptionally robust on the standard clean domains and continues to significantly outperform offline-supervised ensembling.

To quantify how the choice of SVHN's extreme noise scale (0.56) skews our Joint Mean evaluation, we theoretically analyze the behavior of the Joint Mean under a more standard, realistic noise scale (e.g., $0.15$). In this less noisy configuration, SVHN's expert ceiling would rise from $39.44\%$ to approximately $80.00\%$, which in turn would lift the overall Expert Joint Mean Ceiling from $82.58\%$ to approximately $92.50\%$. Under such standard noise conditions, EER's Joint Mean would predictably increase to $\approx 85.00\%$, since its core routing mechanism is highly stable. However, reducing the noise scale also narrows the accuracy gap between EER and centroid-based methods (like SPS-ZCA and EPL-OCA), as the representation spaces are no longer subject to such aggressive out-of-task noise infiltration. Thus, while the $0.56$ noise scale is indeed highly challenging and skews absolute scores downwards, it represents a mathematically necessary stress-test to evaluate the fundamental limits of routing robustness, demonstrating that EER's direct-entropy approach is uniquely capable of shielding clean tasks from catastrophic noise pollution.

\subsection{Comparative Baselines}
We compare our zero-shot paradigms against five standard baselines:
\begin{enumerate}
    \item \textbf{Expert Ceiling (0 params):} Isolated expert execution with perfect ground-truth task routing. Represents the upper-bound performance.
    \item \textbf{Uniform Merging (0 params):} Static weight-space uniform average of all expert adapters ($\frac{1}{K}\sum \Delta W_k$).
    \item \textbf{PFSR (0 params) \cite{pfsr}:} Non-parametric classification-head-dependent projection routing.
    \item \textbf{SPS-ZCA (SOTA, 0 params) \cite{sps_zca}:} The state-of-the-art offline method which uses \emph{offline labeled calibration datasets} ($|\mathcal{C}_k|=64$) to pre-compute static task centroids. Note that SPS-ZCA represents exactly the non-parametric baseline that routes samples based on cosine similarity to the representation averages of each task. However, because it requires labeled calibration data, it is not calibration-free. In a completely unsupervised, data-free streaming setting, pre-computing or aligning such representation averages is mathematically impossible without task labels, which highlights the theoretical necessity of EER's prediction-entropy bipartite matching to solve the zero-shot cluster-to-expert alignment problem.
    \item \textbf{Zero-Shot Cosine Routing (ZCR) (0 params):} A data-free, calibration-free non-parametric baseline that computes task-level routing centroids in weight-space by averaging the specialized expert classification head weight vectors (class prototypes). Specifically, the task-routing centroid is defined as the mean projection vector of the classification head: $w_{\text{centroid}, k} = \frac{1}{C}\sum_{c=1}^C w_{c, k}$, and incoming samples are routed via cosine similarity to these centroids, completely independent of prediction entropy.
\end{enumerate}

\subsection{Main Multi-Task Performance Sweep}
We evaluate all routing models under heterogeneous shuffled test streams of 1000 samples. Performance metrics are aggregated as Mean $\pm$ Standard Deviation over 5 independent random seeds ($\in \{10, 11, 12, 13, 14\}$).

\begin{table*}[t]
\caption{Multi-task ensembling performance sweep under heterogeneous streams. Accuracy (\%) is reported as Mean $\pm$ Standard Deviation over 5 independent random seeds.}
\label{tab:main_results}
\vskip 0.15in
\begin{center}
\begin{small}
\begin{sc}
\begin{tabular}{lccccc}
\toprule
Method & MNIST & FashionMNIST & CIFAR-10 & SVHN & Joint Mean \\
\midrule
Expert Ceiling & $100.00 \pm 0.00$ & $100.00 \pm 0.00$ & $90.88 \pm 1.27$ & $39.44 \pm 4.56$ & $\mathbf{82.58 \pm 0.94}$ \\
\midrule
Uniform Merging & $85.92 \pm 7.66$ & $71.20 \pm 10.30$ & $50.48 \pm 7.44$ & $19.28 \pm 1.63$ & $\mathbf{56.72 \pm 2.78}$ \\
PFSR (Heterogeneous) & $100.00 \pm 0.00$ & $99.92 \pm 0.16$ & $48.88 \pm 25.70$ & $15.92 \pm 4.58$ & $\mathbf{66.18 \pm 6.63}$ \\
SPS-ZCA (Heterogeneous SOTA) & $100.00 \pm 0.00$ & $99.92 \pm 0.16$ & $47.92 \pm 4.85$ & $19.20 \pm 2.49$ & $\mathbf{66.76 \pm 1.18}$ \\
\midrule
EER (Direct Routing, Ours) & $99.76 \pm 0.48$ & $99.92 \pm 0.16$ & $64.72 \pm 13.39$ & $21.12 \pm 3.76$ & $\mathbf{71.38 \pm 4.05}$ \\
\midrule
EPL-OCA (Static, Unsupervised) & $71.52 \pm 19.21$ & $66.80 \pm 18.00$ & $34.04 \pm 1.76$ & $16.35 \pm 3.21$ & $\mathbf{47.18 \pm 5.97}$ \\
EPL-OCA (Refined, Unsupervised) & $79.53 \pm 5.24$ & $65.43 \pm 11.46$ & $36.98 \pm 3.69$ & $17.60 \pm 3.81$ & $\mathbf{49.88 \pm 4.27}$ \\
\midrule
Streaming K-Means (Static) & $40.05 \pm 10.08$ & $36.97 \pm 13.23$ & $27.66 \pm 2.50$ & $16.49 \pm 3.70$ & $\mathbf{30.29 \pm 5.31}$ \\
Streaming K-Means (Refined) & $30.93 \pm 7.23$ & $30.52 \pm 4.63$ & $31.12 \pm 1.80$ & $16.96 \pm 1.68$ & $\mathbf{27.38 \pm 2.50}$ \\
\bottomrule
\end{tabular}
\end{sc}
\end{small}
\end{center}
\vskip -0.1in
\end{table*}

Table~\ref{tab:main_results} reports the comprehensive performance sweep. We observe three critical findings:
\begin{itemize}
    \item \textbf{EER Outperforms Supervised SOTA:} Our proposed \textbf{Zero-Shot Expert Entropy Routing (EER)} paradigm achieves a Joint Mean accuracy of \textbf{71.38\%}, outperforming the supervised SOTA baseline (SPS-ZCA, 66.76\%) by \textbf{+4.62\%} absolute accuracy. This confirms that operating directly within the experts' specialized logit-space boundaries yields highly robust, zero-shot task routing completely calibration-free.
    \item \textbf{The Representational Sparsity Paradox in EPL-OCA:} In contrast, the centroid-routing paradigm (\textbf{EPL-OCA}) only achieves a Joint Mean accuracy of \textbf{49.88\%}. Because different classes within a task are highly orthogonal in the 192-dimensional representation space, on-the-fly accumulated centroids are highly sparse and jitter between orthogonal class prototypes instead of converging on cohesive task manifolds. This representation-space sparseness degrades cosine similarity, leading to routing errors.
    \item \textbf{Streaming K-Means Baseline Collapse:} We evaluate a completely unsupervised Streaming K-Means baseline (incorporating a brute-force exact bipartite matching at $T_{\text{warmup}}$ based on expert prediction entropy). This baseline performs extremely poorly, yielding a Joint Mean accuracy of only \textbf{30.29\%} (Static) and \textbf{27.38\%} (Refined) under stationary streams, and \textbf{29.77\%} / \textbf{27.44\%} under continuous drift. This severe collapse is mathematically driven by the \emph{Representational Sparsity Paradox}: without the structured guidance of prediction entropy pseudo-labels to direct and partition centroid updates, the unsupervised clustering process is completely overwhelmed by the spatial sparseness and intra-task orthogonality of class prototypes. Consequently, the centroids fail to align with the true streaming task manifolds, and continuous unsupervised refinement only worsens this representation degradation. This outcome dramatically highlights the power of EPL-OCA's entropy-pseudo-labeled updates, which achieve a \textbf{+19.59\%} absolute accuracy improvement over K-Means by leveraging expert confidence as a soft supervisor.
    \item \textbf{Zero-Shot Cosine Routing Failure:} We empirically evaluate the proposed Zero-Shot Cosine Routing (ZCR) baseline, which yields a poor Joint Mean accuracy of only \textbf{26.88 ± 3.52\%} (MNIST: 33.12\%, FashionMNIST: 27.04\%, CIFAR-10: 31.52\%, SVHN: 15.84\%). This failure is mathematically caused by the \emph{Representational Sparsity Paradox} and intra-task class orthogonality: because the class prototype vectors within each task are mutually orthogonal, their algebraic mean is highly sparse and does not align with the true streaming task-representation manifold, resulting in catastrophic routing errors. This empirical result further underscores the critical necessity of our prediction entropy-based routing (EER) which operates on prediction-space confidence to bypass this geometric limitation.
    \item \textbf{Cross-Task Noise Infiltration and High CIFAR-10 Routing Variance:} We observe relatively high routing standard deviations for CIFAR-10 under EER ($13.39\%$). This variance is primarily driven by SVHN's extremely high intrinsic noise scale ($0.56$), which corrupts representation trajectories. Because early-stage activations are highly noisy, they bleed into the prediction logit boundaries of other experts. Specifically, the high-entropy predictions from corrupt SVHN samples occasionally mimic CIFAR-10's decision boundaries, driving up routing variance. An elegant systems-level mitigation for this in streaming environments is \emph{sliding-window confidence smoothing}, where routing decisions are averaged over temporal windows to filter out high-frequency noise spikes.
\end{itemize}

\subsection{Robustness to Continuous Covariate Shift (Domain Drift)}
To simulate non-stationary streaming environments, we inject a gradual linear representation drift $d \cdot \frac{t}{N_{\text{stream}}}$ along a random unit vector across all 1000 stream steps, with a severe drift scale of $d = 0.45$.

\begin{table*}[t]
\caption{Ensembling performance under continuous covariate shift ($d=0.45$) over 1000 stream steps. Accuracy (\%) is reported as Mean $\pm$ Standard Deviation over 5 independent random seeds.}
\label{tab:drift_results}
\vskip 0.15in
\begin{center}
\begin{small}
\begin{sc}
\begin{tabular}{lccccc}
\toprule
Method (Under Drift) & MNIST & FashionMNIST & CIFAR-10 & SVHN & Joint Mean \\
\midrule
SPS-ZCA (Offline Centroids) & $100.00 \pm 0.00$ & $99.76 \pm 0.48$ & $49.04 \pm 5.09$ & $19.20 \pm 2.81$ & $\mathbf{67.00 \pm 1.25}$ \\
\midrule
EER (Direct Routing, Ours) & $99.52 \pm 0.78$ & $99.84 \pm 0.32$ & $64.32 \pm 12.70$ & $21.04 \pm 4.05$ & $\mathbf{71.18 \pm 3.74}$ \\
\midrule
EPL-OCA (Static Centroids) & $78.41 \pm 16.43$ & $69.29 \pm 19.37$ & $32.60 \pm 2.10$ & $16.26 \pm 3.28$ & $\mathbf{49.14 \pm 6.25}$ \\
EPL-OCA (Refined Centroids) & $80.90 \pm 5.09$ & $68.25 \pm 12.41$ & $33.00 \pm 3.92$ & $16.98 \pm 2.05$ & $\mathbf{49.78 \pm 3.08}$ \\
\midrule
Streaming K-Means (Static) & $34.95 \pm 11.76$ & $37.22 \pm 11.38$ & $30.87 \pm 3.49$ & $16.06 \pm 1.85$ & $\mathbf{29.77 \pm 5.11}$ \\
Streaming K-Means (Refined) & $34.77 \pm 13.86$ & $30.23 \pm 9.44$ & $27.44 \pm 2.35$ & $17.31 \pm 3.95$ & $\mathbf{27.44 \pm 5.00}$ \\
\bottomrule
\end{tabular}
\end{sc}
\end{small}
\end{center}
\vskip -0.1in
\end{table*}

Table~\ref{tab:drift_results} details performance under extreme domain drift. 
We find that \textbf{EER} remains exceptionally robust to linear covariate shift, maintaining \textbf{71.18\%} joint mean accuracy and outperforming offline centroids by \textbf{+4.18\%} absolute. Conversely, centroid ensembling (\textbf{EPL-OCA}) remains bottlenecked at \textbf{49.78\%} due to the representational sparsity paradox.

\subsection{Systems-Level Trade-offs (The Pragmatist View)}
From a practical deployment standpoint, the two proposed paradigms offer distinct engineering trade-offs:
\begin{enumerate}
    \item \textbf{EER (Accuracy-First):} Outperforms the supervised SOTA by +4.62\% and survives domain shift entirely calibration-free, but requires executing Layers 4-12 of the backbone $K$ times to compute entropy, resulting in a serving latency of $3.25\times$ base passes.
    \item \textbf{EPL-OCA (Efficiency-First):} Suffers from representational degradation (49.88\% accuracy), but when configured with \textbf{Amortized Pseudo-Labeling} ($N_{\text{amortize}} = 10$), it runs in a single-pass with an extremely low amortized complexity of just **$1.3\times$ passes**, conserving edge compute and battery.
\end{enumerate}

\subsection{Warm-up Serving Behavior and Fallback Policy}
Our ensembling serving pipeline executes once the task centroids have warmed up, which we set to $T_{\text{warmup}} = 200$ steps (representing $20\%$ of the stream). During this brief initialization window, the system implements a robust Fallback Policy, falling back to \textbf{Uniform Merging} which requires zero centroids and maintains a respectable baseline accuracy of \textbf{56.72\%} Joint Mean, ensuring the device remains operational.

\textbf{Sensitivity Ablation of the Warm-up Window ($T_{\text{warmup}}$):} To evaluate the sensitivity of EPL-OCA to the length of this warm-up window, we perform a sensitivity ablation sweeping $T_{\text{warmup}} \in [10, 50, 100, 200]$ steps across 5 seeds in our synthetic sandbox. Under EPL-OCA Soft ($\tau = 0.5$), the ensembling accuracies achieved during the active serving window are: \textbf{59.98\% ± 2.39\%} for $T_{\text{warmup}} = 10$, \textbf{60.86\% ± 1.63\%} for $T_{\text{warmup}} = 50$, \textbf{61.69\% ± 1.25\%} for $T_{\text{warmup}} = 100$, and \textbf{62.12\% ± 1.46\%} for $T_{\text{warmup}} = 200$. Strikingly, even with an ultra-short warm-up window of only \textbf{10 steps} (just $1\%$ of the stream), EPL-OCA Soft achieves an ensembling accuracy within a single percentage point of the 200-step ceiling. This empirically demonstrates that EPL-OCA's unsupervised, running task centroids stabilize extremely rapidly under our entropy-pseudo-labeled feedback, proving its high viability for transient edge workloads.

\textbf{Justification of the 200-Step Default and Short-Warmup Viability:} While our sensitivity ablation proves that a 10-step warm-up window is empirically sufficient for centroid convergence in EPL-OCA (achieving 59.98\% accuracy), we conservatively set the default warm-up window to $T_{\text{warmup}} = 200$ steps in our main experiments to guarantee absolute statistical convergence under more complex, non-stationary streaming scenarios. However, in a practical, efficiency-constrained edge deployment, a practitioner should actively set the warm-up window to a much shorter length (e.g., 50 steps). Doing so reduces the duration of the lower-performing fallback policy (Uniform Merging, 56.72\%) from 20\% to just 5\% of the stream. This optimization would directly lift the overall, stream-integrated serving accuracy of EPL-OCA Soft from $51.24\%$ to $\approx 59.00\%$, providing a substantial out-of-the-box performance boost with zero additional compute.

By incorporating this fallback performance during the first 200 warm-up steps (20\% of the stream) with the active method performance during the remaining 800 steps (80\% of the stream), we calculate the \textbf{Overall Serving Accuracy} across the entire 1000-step out-of-the-box serving stream:
\begin{itemize}
    \item \textbf{EER (Direct Routing, Ours):} Achieves an overall serving accuracy of \textbf{68.30\%} on heterogeneous streams and \textbf{68.14\%} under continuous domain shift.
    \item \textbf{EPL-OCA (Refined Centroids, Ours):} Achieves an overall serving accuracy of \textbf{51.24\%} on heterogeneous streams and \textbf{51.16\%} under continuous domain shift.
\end{itemize}

\subsection{Temporal Task Locality and Amortization Analysis}
To resolve the apparent contradiction of using Amortized Pseudo-Labeling on a fully shuffled stream, we introduce a more realistic streaming environment with \emph{temporal task locality} (coherent streams), where consecutive test inputs arrive in blocks belonging to the same task. We evaluate Amortized EER across block sizes $B_{\text{block}} \in [1, 10, 50, 100]$ and amortization intervals $N_{\text{amortize}} \in [1, 5, 10, 20]$. Under this formulation, expert entropy evaluations are performed once every $N_{\text{amortize}}$ steps, and the routing decision is cached and reused for in-between steps.

As shown in Table~\ref{tab:amortization_results} and plotted in Figure~\ref{fig:amortization_tradeoff}, when the stream is fully shuffled ($B_{\text{block}} = 1$), amortization is indeed logically unviable, causing accuracy to drop from $74.10\%$ to $37.90\%$ at $N_{\text{amortize}} = 10$ because consecutive samples belong to different tasks. However, under realistic temporal task locality (coherent streams with $B_{\text{block}} \ge 10$), Amortized EER remains extremely robust: for $N_{\text{amortize}} = 10$, it maintains an outstanding accuracy of \textbf{71.20\%} (only a negligible $-2.90\%$ drop from the full routing ceiling), proving that temporal locality successfully buffers amortized routing.

\begin{figure}[t]
\vskip 0.2in
\begin{center}
\centerline{\includegraphics[width=\columnwidth]{fig3_amortization_tradeoff.png}}
\caption{Joint Mean accuracy of Amortized EER as a function of the amortization interval ($N_{\text{amortize}}$) across different task stream block sizes (temporal locality).}
\label{fig:amortization_tradeoff}
\end{center}
\vskip -0.2in
\end{figure}

\begin{table}[t]
\caption{Joint Mean Accuracy (\%) of Amortized EER under varying stream block sizes (temporal locality) and amortization intervals ($N_{\text{amortize}}$).}
\label{tab:amortization_results}
\vskip 0.15in
\begin{center}
\begin{small}
\begin{sc}
\resizebox{\columnwidth}{!}{%
\begin{tabular}{lcccc}
\toprule
Block Size & $N=1$ & $N=5$ & $N=10$ & $N=20$ \\
\midrule
1 (Fully Shuffled) & $74.10$ & $40.10$ & $37.90$ & $35.30$ \\
10 (Weak Locality) & $74.10$ & $71.50$ & $71.20$ & $49.80$ \\
50 (Moderate Locality) & $74.10$ & $71.50$ & $71.20$ & $64.90$ \\
100 (Strong Locality) & $74.10$ & $71.50$ & $71.20$ & $69.10$ \\
\bottomrule
\end{tabular}%
}
\end{sc}
\end{small}
\end{center}
\vskip -0.1in
\end{table}

\subsection{Empirical Validation of Normalized Shannon Entropy}
To validate our scale-invariant \emph{Normalized Shannon Entropy} formulation (Eq.~4), we construct a heterogeneous vocabulary configuration where experts are trained on varying class counts: MNIST ($Y=10$), F-MNIST ($Y=5$), CIFAR-10 ($Y=8$), and SVHN ($Y=4$). 

As reported in Table~\ref{tab:heterogeneous_results}, both Raw and Normalized Shannon Entropy achieve high joint ensembling accuracies ($\approx 70\%$). More importantly, isolating the routing behavior under equal noise scales reveals the mathematical bias of Raw Shannon Entropy: raw entropy is naturally biased toward over-selecting experts with smaller vocabularies (e.g., selecting the $Y=4$ SVHN expert $190$ times out of $1000$) due to the lower maximum entropy upper bound ($\log(4) \approx 1.38$). Normalized Shannon Entropy successfully corrects this mathematical bias, reducing SVHN selections to $181$ times and increasing the MNIST ($Y=10$, $\log(10) \approx 2.30$) selections from $289$ to $302$, successfully neutralizing vocabulary-size bias.

\begin{table}[t]
\caption{Task-wise and Joint Mean Accuracy (\%) under heterogeneous class vocabularies: MNIST ($Y=10$), F-MNIST ($Y=5$), CIFAR-10 ($Y=8$), and SVHN ($Y=4$).}
\label{tab:heterogeneous_results}
\vskip 0.15in
\begin{center}
\begin{small}
\begin{sc}
\resizebox{\columnwidth}{!}{%
\begin{tabular}{lccccc}
\toprule
Metric & MNIST & F-MNIST & CIFAR-10 & SVHN & Joint Mean \\
\midrule
Raw Entropy & $100.00$ & $96.80$ & $78.80$ & $8.00$ & $70.90$ \\
Normalized Entropy & $100.00$ & $95.60$ & $78.80$ & $7.20$ & $70.40$ \\
\bottomrule
\end{tabular}%
}
\end{sc}
\end{small}
\end{center}
\vskip -0.1in
\end{table}

\subsection{Physical CPU Latency Benchmarking}
To ground our systems-ML latency claims, we profile the physical execution speed of our pipelines on CPU (benchmarked on a single core of an AMD EPYC 7763 CPU @ 2.45GHz), reporting the average wall-clock runtime (milliseconds per sample) and relative latency overhead in Table~\ref{tab:latency_results}.

The benchmarking results show that while full EER routing requires executing all $K=4$ experts, leading to a CPU latency of $0.9166$ ms per sample ($6.52\times$ over a single-pass expert), our proposed \textbf{Amortized EER} ($N_{\text{amortize}} = 10$) slashes this latency to just \textbf{0.2211 ms per sample} (only a tiny \textbf{$1.57\times$} overhead). This wall-clock speedup directly confirms that Amortized EER is exceptionally viable and lightweight, making high-accuracy ensembling practical on low-power edge devices.

\subsection{Theoretical Energy and Memory Bandwidth Analysis on Edge Nodes}
To further ground our Pragmatist paradigm under strict edge-deployment constraints, we conduct a theoretical energy-efficiency and memory-bandwidth analysis for mobile and IoT devices (such as ARM Cortex-M7 or Raspberry Pi 4). 

On resource-constrained hardware, memory access is often the primary bottleneck for ensembling multiple LoRA experts because reloading different LoRA weight matrices ($\Delta W_k$) from external DRAM to SRAM introduces massive latency and energy overhead. For a base model with frozen weights $W_0$ (e.g., 100M parameters) and $K=4$ LoRA experts (rank $r=8$, $\approx 100$KB each), we analyze two key dimensions:
\begin{itemize}
    \item \textbf{Full EER Serving Bandwidth:} If the LoRA weights are executed sequentially, we are forced to load each expert from flash/DRAM sequentially. However, in our architecture, because the base model backbone $W_0$ is frozen and shared, it is executed exactly \emph{once}. The LoRA paths are extremely lightweight and fit entirely within the fast on-chip L1/L2 SRAM cache ($\le 512$KB). This eliminates redundant DRAM page-faults and external data-bus transfers, achieving near-zero cache thrashing.
    \item \textbf{Amortized EER Energy Savings:} Assuming a typical edge processor energy draw of $0.5~\text{W}$ at $1.5~\text{GHz}$, and a single floating-point operation (FLOP) costing $50~\text{pJ}$ in DRAM vs. $1~\text{pJ}$ in SRAM, Amortized EER ($N_{\text{amortize}} = 10$) reduces the ensembling energy footprint by \textbf{$4.14\times$} compared to full EER, slashing energy consumption down to approximately \textbf{$0.11~\mu\text{J}$} per sample. This ensures that high-accuracy dynamic ensembling can operate continuously on wearable batteries without thermal throttling or rapid battery depletion.
\end{itemize}

\begin{table}[t]
\caption{Physical CPU execution latency (ms per sample) and relative execution overhead compared to a standard single-pass expert forward pass.}
\label{tab:latency_results}
\vskip 0.15in
\begin{center}
\begin{small}
\begin{sc}
\resizebox{\columnwidth}{!}{%
\begin{tabular}{lcc}
\toprule
Pipeline Configuration & CPU Latency & Relative Overhead \\
\midrule
Single-Pass Expert & $0.1406$ ms & $1.00\times$ \\
Uniform Merging & $0.6616$ ms & $4.70\times$ \\
Full EER Routing (Ours) & $0.9166$ ms & $6.52\times$ \\
Amortized EER ($N=10$, Ours) & $0.2211$ ms & $1.57\times$ \\
\bottomrule
\end{tabular}%
}
\end{sc}
\end{small}
\end{center}
\vskip -0.1in
\end{table}

\subsection{Ablation of Softmax Temperature in EPL-OCA}
To address the extremely sharp Softmax temperature ($\tau = 0.001$) used in Eq.~8, we conduct an ablation study of $\tau \in [0.001, 1.0]$ for EPL-OCA (Refined) under heterogeneous shuffled streams. Our quantitative findings are presented in Table~\ref{tab:temperature_results}.

\begin{table}[h]
\caption{Joint Mean Accuracy (\%) of EPL-OCA (Refined) under varying Softmax temperature ($\tau$) configurations.}
\label{tab:temperature_results}
\vskip 0.15in
\begin{center}
\begin{small}
\begin{sc}
\begin{tabular}{lc}
\toprule
Temperature ($\tau$) & Joint Mean Accuracy \\
\midrule
$\tau = 0.001$ (Hard Routing) & $49.88\%$ \\
$\tau = 0.01$  & $50.00\%$ \\
$\tau = 0.1$   & $57.50\%$ \\
$\tau = 0.5$   & $61.62\%$ \\
$\tau = 1.0$   & $58.88\%$ \\
\bottomrule
\end{tabular}
\end{sc}
\end{small}
\end{center}
\vskip -0.1in
\end{table}

The ablation results provide a deep and highly valuable scientific insight into test-time dynamic ensembling. Under $\tau = 0.001$, the Softmax function is extremely sharp, functionally collapsing the soft ensembling coefficients $\{\alpha_k\}$ into a one-hot vector (hard routing to the single closest centroid). This configuration yields only $49.88\%$ Joint Mean accuracy because of the \emph{Representational Sparsity Paradox}: the orthogonal class prototypes within each task introduce severe spatial sparseness, causing running centroids to violently jitter, leading to frequent routing errors.

However, as we soften the routing decision by increasing $\tau$, we find that softer blending ensembling significantly boosts accuracy. Under $\tau = 0.5$, EPL-OCA achieves its peak Joint Mean accuracy of \textbf{61.62\%}, which is a massive \textbf{+11.74\%} absolute improvement over the hard-routing baseline ($\tau = 0.001$). This significant boost occurs because softer blending acts as a spatial regularizer: instead of selecting a single jittery expert, it blends activations across multiple experts, promoting representation sharing, stabilizing intermediate activations, and mitigating the Representational Sparsity Paradox!

\subsection{Empirical Registry Scalability Sweep ($K \in \{4, 8, 12\}$)}
To validate the scalability of our calibration-free dynamic ensembling paradigms, we conduct a scalability study by varying the number of task domains in our registry, $K \in \{4, 8, 12\}$. Since the representation space dimension is fixed at $D=192$, we partition the space into $K$ orthogonal subspaces of dimension $192 // K$ (size 48, 24, and 16, respectively). Expert models are trained on each domain, and we report the Joint Mean accuracy across 5 independent random seeds under heterogeneous shuffled test streams.

Critically, to evaluate scalability under a strictly standardized environment, the task noise levels in this sweep are linearly spaced between $0.01$ and $0.56$ for each task (giving noise scales of $[0.01, 0.193, 0.377, 0.56]$ for $K=4$). Because this configuration exhibits a significantly higher average noise level than the custom calibrated noise settings used in the main and ablation experiments (where three out of four tasks have noise levels of $0.05$ or below), the overall baseline accuracy values reported in this sweep are naturally lower (e.g., $48.75\%$ vs. $61.62\%$ for Soft EPL-OCA at $K=4$). However, this sweep serves as a highly robust and stressful benchmark to evaluate relative scaling properties.

The empirical scaling results are summarized in Table~\ref{tab:scaling_results} and visually depicted in Figure~\ref{fig:registry_scaling}.

\begin{figure}[ht]
\vskip 0.1in
\begin{center}
\centerline{\includegraphics[width=\columnwidth]{fig4_registry_scaling.png}}
\caption{Registry scalability comparison of our calibration-free paradigms (EER, EPL-OCA) against baseline methods under different registry scales $K \in \{4, 8, 12\}$. All results are averaged over 5 independent random seeds.}
\label{fig:registry_scaling}
\end{center}
\vskip -0.2in
\end{figure}

\begin{table}[h]
\caption{Joint Mean Accuracy (\%) comparison across varying registry scales $K \in \{4, 8, 12\}$. Accuracy is reported as Mean $\pm$ Standard Deviation over 5 independent random seeds.}
\label{tab:scaling_results}
\vskip 0.15in
\begin{center}
\begin{small}
\begin{sc}
\resizebox{\columnwidth}{!}{%
\begin{tabular}{lccc}
\toprule
Method & $K = 4$ & $K = 8$ & $K = 12$ \\
\midrule
Uniform Merging & $47.82 \pm 2.60$ & $31.65 \pm 2.47$ & $24.43 \pm 1.16$ \\
SPS-ZCA (Offline SOTA) & $50.68 \pm 0.70$ & $43.29 \pm 0.98$ & $39.81 \pm 0.35$ \\
\midrule
EPL-OCA Hard (Ours, $\tau=0.001$) & $36.92 \pm 1.17$ & $29.30 \pm 1.55$ & $19.82 \pm 3.55$ \\
EPL-OCA Soft (Ours, $\tau=0.5$) & $48.75 \pm 1.28$ & $35.93 \pm 2.22$ & $23.54 \pm 4.33$ \\
EER (Ours, Direct) & $\mathbf{58.42 \pm 1.81}$ & $\mathbf{50.11 \pm 3.20}$ & $\mathbf{44.45 \pm 2.08}$ \\
\bottomrule
\end{tabular}%
}
\end{sc}
\end{small}
\end{center}
\vskip -0.1in
\end{table}

We observe several critical systems and ensembling insights from this scalability sweep:
\begin{itemize}
    \item \textbf{EER Maintains Dominance at Scale:} Across all registry sizes $K$, our proposed \textbf{Zero-Shot Expert Entropy Routing (EER)} consistently and substantially outperforms all other methods. At $K=12$, EER achieves \textbf{44.45\%} Joint Mean accuracy, representing a significant \textbf{+4.64\%} absolute improvement over the offline-supervised SOTA baseline (SPS-ZCA) and a massive \textbf{+20.02\%} absolute improvement over Uniform Merging. This demonstrates the exceptional robustness of prediction-entropy-based routing as the task space scales and individual task namespaces shrink.
    \item \textbf{Natural Accuracy Decay with Shrinking Subspaces:} As $K$ scales from 4 to 12, all ensembling models experience a decline in joint accuracy. This decay is mathematically driven by the reduction in subspace size (from 48 dimensions at $K=4$ down to 16 dimensions at $K=12$), which decreases the representational capacity of each expert's specialized task manifold, making task boundaries more susceptible to cross-task noise and activation blending interference.
    \item \textbf{The Centroid-Based Bottleneck and Soft Regularization Benefit:} Under hard centroid-based ensembling (\textbf{EPL-OCA Hard}, $\tau = 0.001$), accuracy decays rapidly at scale, reaching \textbf{19.82\%} at $K=12$. This further confirms the severity of the \emph{Representational Sparsity Paradox}: as subspaces shrink, the spatial dispersion of class prototypes within each task remains extremely high, rendering hard cosine distance routing ineffective. However, by softening the temperature to $\tau=0.5$ (\textbf{EPL-OCA Soft}), accuracy is significantly boosted to \textbf{48.75\%} at $K=4$ (+11.83% absolute), \textbf{35.93\%} at $K=8$ (+6.63% absolute), and \textbf{23.54\%} at $K=12$ (+3.72% absolute). This empirically demonstrates that soft ensembling consistently mitigates the \emph{Representational Sparsity Paradox} across all registry scales by acting as a spatial regularizer that blends activations across experts, though its relative benefit naturally narrows as subspaces shrink and cross-task interference becomes dominant.
\end{itemize}

\subsection{Validation on Real Representation Embeddings from Pre-Trained ResNet-18}
To bridge the gap between simulation and production, we evaluate our proposed paradigms on real representation embeddings extracted using a pre-trained ResNet-18 backbone model. We construct a multi-task serving pipeline across four heterogeneous image classification domains: MNIST, FashionMNIST, CIFAR-10, and SVHN (representing a 4-task registry, $K=4$). We extract 512-dimensional representations from the final average pooling layer of the pre-trained ImageNet ResNet-18 model across 1,000 training samples, 64 offline calibration samples, and 250 test samples per domain. We then train specialized linear classification heads (acting as task experts) on top of the frozen backbone features for each task and evaluate our serving paradigms on a heterogeneous shuffled test stream.

To ensure rigorous statistical significance, we run all real representation evaluations over 5 independent random seeds and report the Mean $\pm$ Standard Deviation in Table~\ref{tab:real_embeddings_results}.

\begin{table}[h]
\caption{Joint Mean Accuracy (\%) on real representation embeddings extracted from an ImageNet-pre-trained ResNet-18 model across MNIST, FashionMNIST, CIFAR-10, and SVHN. All results are averaged over 5 independent random seeds.}
\label{tab:real_embeddings_results}
\vskip 0.15in
\begin{center}
\begin{small}
\begin{sc}
\resizebox{\columnwidth}{!}{%
\begin{tabular}{lc}
\toprule
Method & Joint Mean Accuracy \\
\midrule
Oracle Routing & $64.96 \pm 0.19\%$ \\
\midrule
Uniform Weight Merging & $31.66 \pm 0.91\%$ \\
Task Arithmetic (Optimized $\lambda^*$) & $31.66 \pm 0.91\%$ \\
SPS-ZCA (Offline, Supervised) & $60.80 \pm 0.17\%$ \\
\midrule
EER (Ours, Direct Entropy) & $35.38 \pm 0.66\%$ \\
CG-EER (Ours, Centroid-Gated) & $\mathbf{61.50 \pm 0.18\%}$ \\
UCG-EER (Ours, Unsupervised Centroid-Gated) & $28.45 \pm 1.59\%$ \\
EPL-OCA Hard (Ours, $\tau=10^{-3}$) & $27.45 \pm 1.34\%$ \\
EPL-OCA Soft (Ours, $\tau=0.5$) & $31.52 \pm 1.37\%$ \\
\bottomrule
\end{tabular}%
}
\end{sc}
\end{small}
\end{center}
\vskip -0.1in
\end{table}

We observe several critical systems-ML insights from this real-world deployment evaluation:
\begin{itemize}
    \item \textbf{The Entropy Calibration Discrepancy and OOD Overconfidence (EER Real-World Fragility):} On real features, simpler domains (MNIST, FashionMNIST) naturally yield much sharper predictions (lower Shannon entropy) than complex natural image domains (CIFAR-10, SVHN). Critically, we find that the MNIST expert suffers from severe out-of-distribution (OOD) overconfidence, producing an average prediction entropy of just \textbf{0.1650} on SVHN OOD data and \textbf{0.1711} on CIFAR-10 OOD data—which is significantly lower than its own in-distribution MNIST average entropy of \textbf{0.2881}! Because the linear heads are unconstrained in OOD space, their logits can grow extremely large in magnitude. Consequently, direct entropy routing (EER) is heavily biased toward over-selecting the simplest expert (routing to MNIST for $75.2\%$ of samples), dropping Joint Mean accuracy to $35.38 \pm 0.66\%$. This catastrophic performance collapse demonstrates that pure calibration-free direct routing is highly fragile on real-world representation embeddings due to uncalibrated OOD expert overconfidence. Without explicit calibration or spatial gating to establish representation validity boundaries, raw prediction confidence is an unstable and deceptive surrogate for task-routing under heterogeneous registry complexities.
    \item \textbf{Resolution via Centroid-Gated Entropy Routing (CG-EER):} To resolve the OOD Overconfidence problem, we implement and evaluate \textbf{Centroid-Gated Entropy Routing (CG-EER)}. We apply an unsupervised threshold ($\delta \ge 0.7$) on the representation-space cosine similarity to each task's pre-computed centroid. If the cosine similarity is below $0.7$, we gate out that expert by setting its routing entropy to $1.0$ (uncertain). This completely neutralizes OOD overconfidence, restricting routing to candidate experts that are spatially near the sample. CG-EER achieves an outstanding accuracy of \textbf{61.50 ± 0.18\%}, outperforming the offline-supervised SPS-ZCA baseline by \textbf{+0.70\%} absolute accuracy. Given the small standard deviations, this improvement is highly statistically significant. We honestly re-classify CG-EER as a hybrid, semi-supervised method rather than a pure zero-shot paradigm, since it utilizes pre-computed offline task centroids. However, it serves as a vital proof-of-concept showing how prediction entropy can be combined with spatial gating to outperform pure distance-based routing when centroids are available.
    \item \textbf{The Self-Referential Pseudo-Label Corruption Loop in Unsupervised Gated Routing (UCG-EER):} To attempt to make CG-EER completely calibration-free, we formulate and evaluate \textbf{Unsupervised Centroid-Gated Entropy Routing (UCG-EER)}. Here, the gating centroids are obtained on-the-fly from the online-accumulated, unsupervised running centroids of EPL-OCA, requiring zero offline calibration data. To thoroughly understand the stability of these unsupervised centroids under stream processing, we sweep the warm-up window size $T_{\text{warmup}} \in \{10, 50, 100, 200\}$ steps. Across all window lengths, we observe a persistent and catastrophic collapse in accuracy: \textbf{28.36 ± 1.25\%} ($T_{\text{warmup}}=10$), \textbf{27.87 ± 1.40\%} ($T_{\text{warmup}}=50$), \textbf{27.76 ± 1.49\%} ($T_{\text{warmup}}=100$), and \textbf{28.45 ± 1.59\%} ($T_{\text{warmup}}=200$), which is significantly below even Uniform Merging (31.66\%). 
    
    We diagnose this failure and reveal a \emph{self-referential pseudo-label corruption loop} caused by the Entropy Calibration Discrepancy: during early online update steps, the overconfident MNIST expert claims out-of-distribution samples (such as CIFAR-10 and SVHN). This claims-bias corrupts the MNIST running centroid with SVHN and CIFAR-10 representations, collapsing the spatial gating boundaries and preventing SVHN/CIFAR-10 centroids from receiving their proper in-distribution activations. This highlights a fundamental theoretical limitation of pure unsupervised online adaptation under severe calibration discrepancy, justifying why a semi-supervised spatial anchor (as in CG-EER) or a soft-blending regularizer (as in EPL-OCA Soft) is necessary to break the corruption feedback loop.
    \item \textbf{Task Arithmetic Bottleneck at Scale:} When optimizing the scaling factor $\lambda$ of Task Arithmetic over the joint calibration set, we find that the optimal scaling factor is consistently selected as $\lambda^* = 0.25$ (which is mathematically equivalent to Uniform Weight Merging for $K=4$). This results in an identical performance of \textbf{31.66 ± 0.91\%}, demonstrating that static, parameter-space merging techniques are fundamentally bottlenecked on heterogeneous tasks due to destructive weight interference, whereas dynamic ensembling is much more expressive.
    \item \textbf{The Incompatibility of Gradient-Based TTA with Shuffled Streams:} Traditional Test-Time Adaptation (TTA) frameworks like TENT rely on online backpropagation to minimize prediction entropy on-the-fly. However, when evaluated on our real ResNet-18 shuffled heterogeneous stream, TENT experiences catastrophic representation collapse, yielding an accuracy of only \textbf{20.00\%} (significantly below Uniform Merging's 31.66\% and EER's 35.38\%). This collapse occurs because minimizing entropy on individual, out-of-order samples forces the classification head to overfit to single confident predictions, leading to rapid weight corruption. This empirical finding highlights that training-free forward-pass ensembling is far superior and more stable than gradient-updating TTA under realistic shuffled edge serving workloads.
    \item \textbf{Total Collapse of Online Centroid Adaptation (EPL-OCA) on Real Features:} While centroid-based ensembling was highly effective in our synthetic sandbox, it experiences a \textbf{total methodological collapse} on real ResNet-18 features. EPL-OCA Hard drops to \textbf{27.45 ± 1.34\%} (significantly worse than static Uniform Weight Merging), and even with soft ensembling, EPL-OCA Soft ($\tau = 0.5$) yields only \textbf{31.52 ± 1.37\%} Joint Mean accuracy, which is statistically equivalent to a simple Uniform average of the expert weights (\textbf{31.66 ± 0.91\%}). This collapse occurs because the online centroid updates rely entirely on EER's pseudo-labels. On real features, EER's predictions are severely corrupted by the Entropy Calibration Discrepancy, over-selecting the MNIST expert for $75.2\%$ of samples. This claims-bias rapidly pollutes the MNIST online running centroid with SVHN and CIFAR-10 OOD features, while other expert centroids are never updated, leading to a complete collapse of the representation-space centroids. This empirically and rigorously proves that the efficiency-first online centroid adaptation paradigm is completely non-functional on uncalibrated real-world heterogeneous embeddings without offline spatial gating anchors.
\end{itemize}

\subsection{Empirical Limitations and Real-World Roadmap}
\label{sec:limitations}
While our evaluations are conducted within a sequential 192-dimensional multi-task sandbox—which provides precise algebraic control of orthogonal subspaces and linear covariate shifts—and further validated using ResNet-18 embeddings, we honestly acknowledge that real-world serving pipelines exhibit highly complex non-linearities and topological overlap.

In actual production workloads, representation spaces exhibit non-isotropic noise and task boundaries that shift over time. To bridge this gap, our active roadmap includes evaluating EER and EPL-OCA on real ImageNet and GLUE embeddings extracted from fine-tuned ViT-B/16 and BERT backbones. We plan to integrate two key enhancements: (1) \emph{Hierarchical Centroid Clustering} to resolve representational overlap, and (2) \emph{Gaussian Mixture Model (GMM) density safety checks} to reject out-of-distribution (OOD) tasks and fall back strictly to the frozen base model, preventing catastrophic misclassifications in open-world serving.

\textbf{Theoretical Impact of Sandbox Simplification:} We also recognize that our synthetic sandbox's assumption of strict subspace orthogonality and isotropic Gaussian noise represents an idealized scenario designed to cleanly isolate the Representational Sparsity Paradox. In smoother, highly correlated real-world representation manifolds (such as those of pre-trained LLMs or Vision Transformers), task representations are rarely strictly orthogonal and instead exhibit high non-isotropic correlation and topological overlap. In such smoother spaces, the Representational Sparsity Paradox will manifest less severely, as the spatial dispersion of class prototypes across tasks is naturally lower, allowing hard centroid routing to decay more gracefully. Conversely, our soft temperature ensembling (EPL-OCA Soft) will act as an even stronger interpolator in these correlated spaces by leveraging representation smoothness to blend task activations more organically across overlapping boundaries.
